\section{Related Work}
\label{sec:related-work}

\paragraph{Moment retrieval models.}
Recent work on natural-language video localization has explored both span-prediction and proposal-style formulations. 2D-TAN represents candidate moments on a two-dimensional start-end map so the model can capture relations between adjacent temporal spans \cite{zhang2019twodtan}. VSLNet instead casts localization as a span prediction problem inspired by question answering and introduces query-guided highlighting to focus the search region \cite{zhang2020vslnet}. The \dataset benchmark pushed the task toward a more realistic multimodal setting by pairing TV clips with subtitles and by annotating whether descriptions depend on visual evidence, subtitle evidence, or both \cite{lei2020tvr}. The \xml model introduced with \dataset is especially relevant here because it provides a strong cross-modal baseline for multimodal moment retrieval.

\paragraph{Retrieval-oriented and hierarchical search.}
Not all systems treat exhaustive multimodal scoring as the first step. Escorcia et al.\ study moment retrieval over large video collections and show that an indexable retrieval stage followed by finer alignment can make the search substantially more scalable \cite{escorcia2019finding}. That retrieval-first perspective is close to the final design in this project. We use lexical retrieval, dense text retrieval, and fusion as inexpensive subtitle-based probes, then use dense retrieval as a routing stage before \xml. This keeps the expensive cross-modal model in the pipeline, but asks it to score a smaller candidate set.

\paragraph{Efficient video-language systems.}
Systems work has increasingly treated multimodal inference cost as a core design concern rather than a secondary implementation detail. Zelda shows that vision-language models can support expressive natural-language video analytics, while still requiring careful ranking and pruning to keep top-$k$ search practical \cite{romero2023zelda}. Deja Vu pushes efficiency further by reusing computation across adjacent frames in ViT-based video-language models and pairing that reuse with memory-compute compaction so that algorithmic savings translate into real GPU speedups \cite{hwang2025dejavu}. Our work is complementary to these directions. Rather than modifying the visual backbone or assuming access to raw frames, we operate in a fixed-feature \dataset setting and reduce validation cost through routing, caching, and a separate temporal-compaction ablation.

Taken together, these lines of work motivate the central question of the report: if \xml is the model that localizes well, how much of the search around it can be made cheap without losing the moment?
